% GENERATED FILE. Edit canonical lessons, then run npm run generate:books.
% canonical-source-hash: fnv1a64:3c51cc2e1a56301a
% canonical-lessons: TA-C444-cintu, TA-C444-citaru, TA-C444-teri, TA-C444-karai-stain, TA-C444-oluku

\chapter{Spills and Stains}
\label{ch:ta-spills-and-stains}
\hlchaptermodality{444}

\begin{chapteropening}
\textbf{By the end of this chapter:} \emph{I can talk about spilling, scattering and splashing, stains, and a leaking roof.}

Complete the last lesson of chapter 444: I can talk about spilling, scattering and splashing, stains, and a leaking roof.
\end{chapteropening}

\section[\texorpdfstring{\ta{சிந்து} (cintu)}{சிந்து (cintu)}]{\ta{சிந்து} (cintu) — to spill}

\label{lesson:TA-C444-cintu}



\begin{quote}
Before the new one: say the Tamil for a hole, then the Tamil for to scratch.
\end{quote}

\subsection*{You'll want to know: \ta{சிந்து}}
\textbf{\ta{சிந்து}} — \emph{cintu} — \textquotedblleft{}to spill\textquotedblright{}.

\begin{quote}
\ta{நான்} \ta{தரையில்} \ta{காபியைச்} \ta{சிந்தினேன்}.
\end{quote}

\begin{quote}
\emph{nāṉ taraiyil kāpiyaic cintiṉēṉ.}
\end{quote}

\begin{quote}
I spilled coffee on the floor.
\end{quote}

Past (I …): \ta{சிந்தினேன்}, \emph{cintiṉēṉ}.

Also of tears: \ta{கண்ணீர்} \ta{சிந்து}, to shed tears.

\begin{cousinweb}
\textbf{\ta{சிந்து}} (\emph{cintu}) is native Dravidian (DEDR 1546), \textquotedblleft{}to be spilled, to trickle; to scatter\textquotedblright{}.
\end{cousinweb}

\begin{grammarlens}[title={What you've built}]
One more word for getting around.
\end{grammarlens}

\subsection*{Your turn}
\begin{itemize}
\raggedright
  \item \emph{Say it:} \emph{cintu}
  \item \emph{Say it:} \emph{cintu}, once more
  \item \emph{Say it:} \emph{nāṉ taraiyil kāpiyaic cintiṉēṉ.} — I spilled coffee on the floor.
  \item \emph{Recall:} say the Tamil for a hole, then the Tamil for to scratch, then say \emph{cintu} again
\end{itemize}

\begin{tcolorbox}[breakable,colback=teal!4,colframe=teal!35!black,title={Before you move on}]
What does \ta{சிந்து} mean? (\textquotedblleft{}To spill\textquotedblright{}.) Say it once more.
\end{tcolorbox}
\section[\texorpdfstring{\ta{சிதறு} (citaṟu)}{சிதறு (citaṟu)}]{\ta{சிதறு} (citaṟu) — to scatter}

\label{lesson:TA-C444-citaru}



\begin{quote}
Before the new one: say the Tamil for to scratch, then the Tamil for to spill.
\end{quote}

\subsection*{You'll want to know: \ta{சிதறு}}
\textbf{\ta{சிதறு}} — \emph{citaṟu} — \textquotedblleft{}to scatter\textquotedblright{}.

\begin{quote}
\ta{பை} \ta{கிழிந்ததால்} \ta{அரிசி} \ta{தரையில்} \ta{சிதறியது}.
\end{quote}

\begin{quote}
\emph{pai kiḻintatāl arici taraiyil citaṟiyatu.}
\end{quote}

\begin{quote}
Because the bag tore, the rice scattered on the floor.
\end{quote}

Past (it …): \ta{சிதறியது}, \emph{citaṟiyatu}.

Of a crowd: \ta{மக்கள்} \ta{சிதறி} \ta{ஓடினார்கள்}, the people scattered.

\begin{cousinweb}
\textbf{\ta{சிதறு}} (\emph{citaṟu}) is native Dravidian (DEDR 1546), \textquotedblleft{}to disperse, to scatter, to be shed\textquotedblright{}.
\end{cousinweb}

\begin{grammarlens}[title={What you've built}]
One more word for getting around.
\end{grammarlens}

\subsection*{Your turn}
\begin{itemize}
\raggedright
  \item \emph{Say it:} \emph{citaṟu}
  \item \emph{Say it:} \emph{citaṟu}, once more
  \item \emph{Say it:} \emph{pai kiḻintatāl arici taraiyil citaṟiyatu.} — Because the bag tore, the rice scattered on the floor.
  \item \emph{Recall:} say the Tamil for to scratch, then the Tamil for to spill, then say \emph{citaṟu} again
\end{itemize}

\begin{tcolorbox}[breakable,colback=teal!4,colframe=teal!35!black,title={Before you move on}]
What does \ta{சிதறு} mean? (\textquotedblleft{}To scatter\textquotedblright{}.) Say it once more.
\end{tcolorbox}
\section[\texorpdfstring{\ta{தெறி} (teṟi)}{தெறி (teṟi)}]{\ta{தெறி} (teṟi) — to splash; to spatter}

\label{lesson:TA-C444-teri}



\begin{quote}
Before the new one: say the Tamil for to spill, then the Tamil for to scatter.
\end{quote}

\subsection*{You'll want to know: \ta{தெறி}}
\textbf{\ta{தெறி}} — \emph{teṟi} — \textquotedblleft{}to splash; to spatter\textquotedblright{}.

\begin{quote}
\ta{வண்டி} \ta{போனபோது} \ta{சேறு} \ta{என்} \ta{மேல்} \ta{தெறித்தது}.
\end{quote}

\begin{quote}
\emph{vaṇṭi pōṉapōtu cēṟu eṉ mēl teṟittatu.}
\end{quote}

\begin{quote}
When the vehicle went past, mud splashed on me.
\end{quote}

Past (it …): \ta{தெறித்தது}, \emph{teṟittatu}.

Where it lands takes \ta{மேல்}: \ta{என்} \ta{மேல்} \ta{தெறித்தது}.

\begin{cousinweb}
\textbf{\ta{தெறி}} (\emph{teṟi}) is native Dravidian (DEDR 3438), \textquotedblleft{}to strike and fly off, to splash\textquotedblright{}.
\end{cousinweb}

\begin{grammarlens}[title={What you've built}]
One more word for getting around.
\end{grammarlens}

\subsection*{Your turn}
\begin{itemize}
\raggedright
  \item \emph{Say it:} \emph{teṟi}
  \item \emph{Say it:} \emph{teṟi}, once more
  \item \emph{Say it:} \emph{vaṇṭi pōṉapōtu cēṟu eṉ mēl teṟittatu.} — When the vehicle went past, mud splashed on me.
  \item \emph{Recall:} say the Tamil for to spill, then the Tamil for to scatter, then say \emph{teṟi} again
\end{itemize}

\begin{tcolorbox}[breakable,colback=teal!4,colframe=teal!35!black,title={Before you move on}]
What does \ta{தெறி} mean? (\textquotedblleft{}To splash; to spatter\textquotedblright{}.) Say it once more.
\end{tcolorbox}
\section[\texorpdfstring{\ta{கறை} (kaṟai)}{கறை (kaṟai)}]{\ta{கறை} (kaṟai) — a stain}

\label{lesson:TA-C444-karai-stain}



\begin{quote}
Before the new one: say the Tamil for to scatter, then the Tamil for to splash.
\end{quote}

\subsection*{You'll want to know: \ta{கறை}}
\textbf{\ta{கறை}} — \emph{kaṟai} — \textquotedblleft{}a stain\textquotedblright{}.

\begin{quote}
\ta{சட்டையில்} \ta{காபிக்} \ta{கறை} \ta{இருக்கிறது}.
\end{quote}

\begin{quote}
\emph{caṭṭaiyil kāpik kaṟai irukkiṟatu.}
\end{quote}

\begin{quote}
There is a coffee stain on the shirt.
\end{quote}

\ta{கறை} \ta{படி} is to get stained.

\begin{cousinweb}
\textbf{\ta{கறை}} (\emph{kaṟai}) is native Dravidian (DEDR 1395), \textquotedblleft{}a spot, a stain, a blemish\textquotedblright{}.
\end{cousinweb}

\begin{grammarlens}[title={What you've built}]
One more word for getting around.
\end{grammarlens}

\subsection*{Your turn}
\begin{itemize}
\raggedright
  \item \emph{Say it:} \emph{kaṟai}
  \item \emph{Say it:} \emph{kaṟai}, once more
  \item \emph{Say it:} \emph{caṭṭaiyil kāpik kaṟai irukkiṟatu.} — There is a coffee stain on the shirt.
  \item \emph{Recall:} say the Tamil for to scatter, then the Tamil for to splash, then say \emph{kaṟai} again
\end{itemize}

\begin{tcolorbox}[breakable,colback=teal!4,colframe=teal!35!black,title={Before you move on}]
What does \ta{கறை} mean? (\textquotedblleft{}A stain\textquotedblright{}.) Say it once more.
\end{tcolorbox}
\section[\texorpdfstring{\ta{ஒழுகு} (oḻuku)}{ஒழுகு (oḻuku)}]{\ta{ஒழுகு} (oḻuku) — to drip; to leak}

\label{lesson:TA-C444-oluku}



\begin{quote}
Before the new one: say the Tamil for to splash, then the Tamil for a stain.
\end{quote}

\subsection*{You'll want to know: \ta{ஒழுகு}}
\textbf{\ta{ஒழுகு}} — \emph{oḻuku} — \textquotedblleft{}to drip; to leak\textquotedblright{}.

\begin{quote}
\ta{மழையில்} \ta{கூரை} \ta{ஒழுகியது}.
\end{quote}

\begin{quote}
\emph{maḻaiyil kūrai oḻukiyatu.}
\end{quote}

\begin{quote}
The roof leaked in the rain.
\end{quote}

Past (it …): \ta{ஒழுகியது}, \emph{oḻukiyatu}.

Of a tap: \ta{குழாய்} \ta{ஒழுகுகிறது}, the tap is dripping.

\begin{cousinweb}
\textbf{\ta{ஒழுகு}} (\emph{oḻuku}) is native Dravidian (DEDR 1010), \textquotedblleft{}to flow, to leak, to trickle down\textquotedblright{}.
\end{cousinweb}

\begin{grammarlens}[title={What you've built}]
One more word for getting around.
\end{grammarlens}

\subsection*{Your turn}
\begin{itemize}
\raggedright
  \item \emph{Say it:} \emph{oḻuku}
  \item \emph{Say it:} \emph{oḻuku}, once more
  \item \emph{Say it:} \emph{maḻaiyil kūrai oḻukiyatu.} — The roof leaked in the rain.
  \item \emph{Recall:} say the Tamil for to splash, then the Tamil for a stain, then say \emph{oḻuku} again
\end{itemize}

\begin{tcolorbox}[breakable,colback=teal!4,colframe=teal!35!black,title={Before you move on}]
What does \ta{ஒழுகு} mean? (\textquotedblleft{}To drip; to leak\textquotedblright{}.) Say it once more.
\end{tcolorbox}
